\documentclass[12pt, a4paper, openany]{report}

\def\VersionRapport{1.0}

\usepackage[utf8]{inputenc} % un package
\usepackage[T1]{fontenc}      % un second package
\usepackage[francais,english]{babel}  % un troisième package
\usepackage{layout}
\usepackage[top=2.7cm, bottom=2.5cm, left=3.5cm, right=3cm]{geometry}
\usepackage{setspace}

\frenchbsetup{StandardLists=true} % à inclure si on utilise \usepackage[french]{babel}
\usepackage{enumitem}
\usepackage{amssymb}

\usepackage{color}
\usepackage{listings}
\definecolor{dkgreen}{rgb}{0,0.6,0}
\definecolor{gray}{rgb}{0.5,0.5,0.5}
\definecolor{mauve}{rgb}{0.58,0,0.82}

\lstset{frame=tb,
  language=Java,
  aboveskip=3mm,
  belowskip=3mm,
  showstringspaces=false,
  columns=flexible,
  basicstyle={\small\ttfamily},
  numbers=none,
  numberstyle=\tiny\color{gray},
  keywordstyle=\color{blue},
  commentstyle=\color{dkgreen},
  stringstyle=\color{mauve},
  breaklines=true,
  breakatwhitespace=true,
  tabsize=3
}

\usepackage{multirow} % pour les tableaux
\usepackage[table]{xcolor} % pour les tableaux

\usepackage{verbatim}
\usepackage{moreverb}
\usepackage{url}
\usepackage{pst-all}
\usepackage{eso-pic,graphicx}
\usepackage{caption} 
\usepackage[colorlinks=true,urlcolor=blue,linkcolor=red]{hyperref}
\usepackage{array}
\usepackage[toc,page]{appendix}
\usepackage[off]{auto-pst-pdf}
\usepackage{hyperref} % pour le sommaire table des matières
\AddThinSpaceBeforeFootnotes % à insérer si on utilise \usepackage[french]{babel}
\FrenchFootnotes % à insérer si on utilise \usepackage[french]{babel}
\usepackage{fancyhdr}
\pagestyle{headings}

\renewcommand{\appendixpagename}{Annexes}
\renewcommand{\appendixtocname}{Annexes}

\title{Theme: Performance & Robustesse}
\author{Mehenna \bsc{Rebout}}
\date{2018-2019}



%new
\newcommand{\HRule}{\rule{\linewidth}{0.5mm}}


\begin{document}

%\selectlanguage{francais}
\pagenumbering{arabic} 

\makeatletter
  \begin{titlepage}
  

  \begin{sffamily}
   \begin{center}

    % Upper part of the page. The '~' is needed because \\
    % only works if a paragraph has started.
    \includegraphics[scale=0.5]{Logo_UT3.jpg}~\\[1.5cm]

    \textsc{\LARGE UFR EEA }\\[2cm]

    \textsc{\Large Rapport Perfo et Robust}\\[1.5cm]

    % Title
    \HRule \\[0.4cm]
    { \huge \bfseries Archi des Pics\\[0.4cm] }

    \HRule \\[2cm]
    \includegraphics[scale=0.3]{Logo_UT3.jpg}
    \\[2cm]

    % Author and supervisor
    \begin{minipage}{0.4\textwidth}
      \begin{flushleft} \large
         \textsc{Ali Kherbiche}\\
        Promo \\
      \end{flushleft}
    \end{minipage}
    \begin{minipage}{0.4\textwidth}
      \begin{flushright} \large
        \emph{Tuteur :}  \textsc{M. GOUESBAUT}\\
        \emph{Responsable de la Formation:} \textsc{M. BIDON}
      \end{flushright}
    \end{minipage}

    \vfill

    % Bottom of the page
    {\large Nov 2018}

  \end{center}
  \end{sffamily}      
          
  \end{titlepage}
  
\makeatother



% *********************** Remerciements *****************
\chapter*{Remerciements}
 \addcontentsline{toc}{chapter}{Remerciements}

  Je tiens à exprimer ma profonde gratitude à mon promoteur, monsieur XXXX pour m'avoir encadré et guidé tout au long de l'année scolaire et mon mémoire, pour ses conseils judicieux et minutieusement prodigués.\\
  
  Aussi je tiens à lui reconnaître le temps précieux qu’il m'a consacré \\
  
   Que les membres du jury trouvent ici mes remerciements les plus vifs pour avoir accepté d’honorer par leur jugement mon travail.\\
   
   Mes sincères sentiments vont à tous ceux qui, de près ou de loin, ont contribué à la réalisation des mes études supérieurs, en particulier ma chères familles et mes amis (es).\\
   
   
% *********************** Dédicaces *****************
\chapter*{Dédicaces}
 \addcontentsline{toc}{chapter}{Dédicaces}
 \begin{center}
	\begin{Large} \textbf{À mes très chers parents} \end{Large}\\ 
	\begin{Large} \textbf{À mes frères} \end{Large}\\
	\begin{Large} \textbf{À toute ma famille} \end{Large}\\
	\begin{Large} \textbf{À tous mes amis (es)} \end{Large} \\ 
	\begin{Large} \textbf{À toute la promo} \end{Large}\\
 \end{center}
 \begin{flushright} \begin{Large} \textit{Ali Kherbiche} \end{Large} \end{flushright}

%*********************** somaire **************
\tableofcontents
%*********************** listes des figures **************
\listoffigures
%*********************** listes des tableaux **************
\listoftables



%*********************** INTRODUCTION **************
\chapter*{Introduction}
 \addcontentsline{toc}{chapter}{Introduction}
 
  Une des raisons de construire des logiciels est pour une utilité utilisateurs.bla bla
                                                      



%*********************** Problématique **************
\chapter*{Problématique}
 \addcontentsline{toc}{chapter}{Problématique}
   Depuis l’apparition de la nécessité de. \\
   
   Les systèmes .... bla bla... ajoutant à cela quelques problèmes connus:
   
    \begin{itemize}[label=$\square$]
      \item  Les applications.
      \item  Les problèmes de. 
      \item  L'ajout de.
      \item  Les applications.
  \end{itemize}
  
   Malheureusement beaucoup d'entreprises ont bla bla. \\
   
   De nos jours....
   
   



%*********************** État de l'art **************
\chapter*{État de l'art}
\addcontentsline{toc}{chapter}{État de l'art}
 \chaptermark{État de l'art}

 Les années 1960 ont été le début du bla bla .... \\
 
 L'apparition de bla bla ....\\
 
 L'architecture bla bla...\\
 
 L'architecture....
 
 Maintenant, bla bla....
 
 Les bla bla....
 



%*********************** Performance et Robustesse **************
\chapter{Performance et Robustesse}
 La Performance c’est tout d’abord un blaba.\\
 
 La Robustesse c’est tout d’abord un blaba.\\ \\
 % captures d’écrans 
 \begin{center}
   \includegraphics[scale=0.5]{Logo_UT3.jpg}
   \captionof{figure}{\textit{Exemple d'architecture}}
   \label{fig1}
 \end{center}
 
 La figure au-dessus a ... bla bla ....\\
 
 Les bla bla ....\\
 
 Néanmoins, cette structure n’est jamais stable au fil du temps, ... bla bla.\\
  
 En court, il y a une différence entre ... bla bla ..., est fort envisageable qu’elle subira des changements dans le temps.
 
 \paragraph{Définition:}
  Bla bla ....

 \section{Section une}
  ggggggggggggggg

\end{document}
